% Part: first-order-logic
% Chapter: model-theory
% Section: partial-iso

\documentclass[../../../include/open-logic-section]{subfiles}

\begin{document}

\olfileid{mod}{bas}{dlo}
\section{Dichte lineaire ordeningen}

\begin{defn}
  Een \emph{dichte lineaire ordening zonder eindpunten} is een !!a{structure}
  $\Struct{M}$ voor de !!{language} die één tweeplaatsig
  !!{predicate}~$<$ bevat en aan de volgende zinnen voldoet:
  \begin{enumerate}
  \item $\lforall[x][\lnot x < x]$;
  \item $\lforall[x][\lforall[y][\lforall[z][(x < y \lif (y < z \lif x
    <z ))]]]$;
  \item $\lforall[x][\lforall[y][(x< y \lor \eq[x][y] \lor y < x)]]$;
  \item $\lforall[x][\lexists[y][x < y]]$;
  \item $\lforall[x][\lexists[y][y < x]]$;
  \item $\lforall[x][\lforall[y][(x < y \lif \lexists[z][(x < z \land
        z < y))]]]$.
 \end{enumerate}
\end{defn}

\begin{thm}\ollabel{thm:cantorQ}
  Elk tweetal !!{enumerable} dichte lineaire ordeningen zonder
  eindpunten is isomorf.
\end{thm}

\begin{proof}
  Zij $\Struct{M_1}$ en $\Struct{M_2}$ !!{enumerable} dichte lineaire
  ordeningen zonder eindpunten, met ${<_1} = \Assign{<}{M_1}$ en ${<_2} =
  \Assign{<}{M_2}$, en zij $\PIso{I}$ de verzameling van alle partiële
  isomorfismen tussen beide. $\PIso{I}$ is niet leeg, want in elk geval
  $\emptyset \in \PIso{I}$. We tonen aan dat $\PIso{I}$ aan de
  heen-en-weereigenschap voldoet. Dan geldt $\Struct{M_1} \iso[p] \Struct{M_2}$
  en volgt de stelling uit \olref[pis]{thm:p-isom1}.

  Om aan te tonen dat $\PIso{I}$ aan de heen-eigenschap voldoet, zij $p \in
  \PIso{I}$ en zij $p(a_i) = b_i$ voor $i = 1$, \dots,~$n$. Neem
  zonder verlies van algemeenheid aan dat $a_1 <_1 a_2 <_1 \cdots <_1
  a_n$. Gegeven $a \in \Domain{M_1}$ vinden we als volgt een
  $b \in \Domain{M_2}$:
  \begin{enumerate}
  \item als $a <_1 a_1$, kies $b \in \Domain{M_2}$ zo dat $b <_2
    b_1$;
  \item als $a_n <_1 a$, kies $b \in \Domain{M_2}$ zo dat $b_n <_2 b$;
 \item als $a_i <_1 a <_1 a_{i+1}$ voor zekere $i$, kies dan $b \in
   \Domain{M_2}$ zo dat $b_i <_2 b <_2 b_{i+1}$.
  \end{enumerate}
  Er kan steeds een $b$ met de gewenste eigenschap worden gevonden, want
  $\Struct{M_2}$ is een dichte lineaire ordening zonder eindpunten. Definieer
  $q = p \cup \{ \langle a, b \rangle \}$, zodat $q \in \PIso{I}$ de
  gewenste uitbreiding van $p$ is. Hiermee is de heen-eigenschap
  bewezen. De weer-eigenschap gaat analoog. Dus $\Struct{M_1} \iso[p]
  \Struct{M_2}$; volgens \olref[pis]{thm:p-isom1} geldt $\Struct{M_1} \iso
  \Struct{M_2}$.
\end{proof}

\begin{prob}
  Voltooi het bewijs van \olref[mod][bas][dlo]{thm:cantorQ} door
  te verifiëren dat $\PIso{I}$ aan de weer-eigenschap voldoet.
\end{prob}

\begin{rem}
  Zij $\Struct{S}$ een willekeurige !!{enumerable} dichte lineaire ordening zonder
  eindpunten. Dan geldt (volgens \olref{thm:cantorQ}) $\Struct{S} \iso
  \Struct{Q}$, waarbij $\Struct{Q} = (\Rat, <)$ de !!{enumerable}
  dichte lineaire ordening is waarvan de verzameling $\Rat$ van rationale getallen
  het domein vormt. Beschouw nu opnieuw de !!{structure}~$\Struct{R} =
  (\Real, <)$ uit \olref[thm]{remark:R}. We hebben gezien dat er
  !!a{enumerable} !!{structure}~$\Struct{S}$ bestaat waarvoor $\Struct{R}
  \elemequiv \Struct{S}$. Maar $\Struct{S}$ is !!a{enumerable} dichte
  lineaire ordening zonder eindpunten en is dus isomorf (en bijgevolg
  elementair equivalent) met de !!{structure}~$\Struct{Q}$. Uit de
  transitiviteit van elementaire equivalentie volgt $\Struct{R} \elemequiv
  \Struct{Q}$. (Dit had ook rechtstreeks kunnen worden aangetoond door met
  hetzelfde heen-en-weerargument $\Struct{R} \iso[p] \Struct{Q}$ vast te
  stellen.)
\end{rem}
\end{document}
